% ============================================================================
%  El anillo de firmas
%  Módulo 08 · Privacidad — ChainNotes
%
%  Versión LaTeX del apunte web temas/anillo-de-firmas.html
%  Todas las figuras se generan con TikZ: no hay imágenes externas.
%  Los ejemplos numéricos son los mismos que produce assets/js/anillo-de-firmas.js
%  con p = 23, q = 11, g = 4.
%
%  Compilar:  pdflatex anillo-de-firmas.tex   (dos pasadas, por el índice)
%             o bien:  latexmk -pdf anillo-de-firmas.tex
% ============================================================================

\documentclass[11pt,a4paper]{article}

% --- idioma y codificación -------------------------------------------------
\usepackage[T1]{fontenc}
\usepackage[utf8]{inputenc}
\usepackage[spanish,es-noshorthands,es-tabla]{babel}
\usepackage{lmodern}
\usepackage{microtype}

% --- página ----------------------------------------------------------------
\usepackage[a4paper,margin=2.6cm,headheight=14pt]{geometry}
\usepackage{fancyhdr}
\usepackage{enumitem}
\usepackage{booktabs}
\usepackage{tabularx}
\usepackage{array}
\usepackage{caption}
\usepackage{multicol}

% --- matemáticas -----------------------------------------------------------
\usepackage{amsmath,amssymb,amsthm}
\usepackage{mathtools}
\usepackage{pifont}   % \ding{172} y \ding{173}: las condiciones (1) y (2) del eslabón

% --- gráficas --------------------------------------------------------------
\usepackage{tikz}
\usetikzlibrary{positioning,arrows.meta,calc,fit,backgrounds,shapes.geometric,
                decorations.pathreplacing}

% --- enlaces y cajas (hyperref antes que tcolorbox) ------------------------
\usepackage[colorlinks=true,linkcolor=violetdeep,urlcolor=violetdeep,
            citecolor=violetdeep,pdftitle={El anillo de firmas},
            pdfauthor={ChainNotes}]{hyperref}
\usepackage[most]{tcolorbox}

% ---------------------------------------------------------------------------
%  Paleta (heredada de assets/css/styles.css)
% ---------------------------------------------------------------------------
\definecolor{violetc}{HTML}{7A68F5}
\definecolor{violetdeep}{HTML}{5B46E0}
\definecolor{mintc}{HTML}{1FA88A}
\definecolor{coralc}{HTML}{E04E5C}
\definecolor{amberc}{HTML}{C98A1F}
\definecolor{skyc}{HTML}{1E88C7}
\definecolor{inkc}{HTML}{282344}

\hypersetup{linkcolor=violetdeep}

% ---------------------------------------------------------------------------
%  Estilo de párrafo y encabezados
% ---------------------------------------------------------------------------
\setlength{\parindent}{0pt}
\setlength{\parskip}{0.55em}
\captionsetup{font=small,labelfont={bf,color=violetdeep}}

\renewcommand{\topfraction}{0.9}
\renewcommand{\bottomfraction}{0.8}
\renewcommand{\textfraction}{0.07}
\renewcommand{\floatpagefraction}{0.75}
\setcounter{topnumber}{3}
\setcounter{totalnumber}{4}

\pagestyle{fancy}
\fancyhf{}
\fancyhead[L]{\small\color{inkc}\textsc{El anillo de firmas}}
\fancyhead[R]{\small\color{inkc}\thepage}
\renewcommand{\headrulewidth}{0.4pt}

% ---------------------------------------------------------------------------
%  Entornos de teoremas
% ---------------------------------------------------------------------------
\theoremstyle{definition}
\newtheorem{definicion}{Definición}
\theoremstyle{plain}
\newtheorem{proposicion}{Proposición}
\theoremstyle{remark}
\newtheorem*{observacion}{Observación}

% ---------------------------------------------------------------------------
%  Cajas de color
% ---------------------------------------------------------------------------
\newtcolorbox{nota}[1]{enhanced,breakable,colback=skyc!6,colframe=skyc!75!black,
  coltitle=white,fonttitle=\bfseries\small,title={#1},
  left=7pt,right=7pt,top=5pt,bottom=5pt,boxrule=0.7pt,arc=2.5pt,
  before upper={\setlength{\parskip}{0.5em}}}

\newtcolorbox{caso}[1]{enhanced,breakable,colback=amberc!8,colframe=amberc!85!black,
  coltitle=white,fonttitle=\bfseries\small,title={#1},
  left=7pt,right=7pt,top=5pt,bottom=5pt,boxrule=0.7pt,arc=2.5pt,
  before upper={\setlength{\parskip}{0.5em}}}

\newtcolorbox{aviso}[1]{enhanced,breakable,colback=coralc!7,colframe=coralc!85!black,
  coltitle=white,fonttitle=\bfseries\small,title={#1},
  left=7pt,right=7pt,top=5pt,bottom=5pt,boxrule=0.7pt,arc=2.5pt,
  before upper={\setlength{\parskip}{0.5em}}}

\newtcolorbox{clave}[1]{enhanced,breakable,colback=mintc!7,colframe=mintc!85!black,
  coltitle=white,fonttitle=\bfseries\small,title={#1},
  left=7pt,right=7pt,top=5pt,bottom=5pt,boxrule=0.7pt,arc=2.5pt,
  before upper={\setlength{\parskip}{0.5em}}}

\newtcolorbox{laboratorio}[1]{enhanced,breakable,colback=violetc!6,colframe=violetc!80!black,
  coltitle=white,fonttitle=\bfseries\small,title={#1},
  left=7pt,right=7pt,top=5pt,bottom=5pt,boxrule=0.7pt,arc=2.5pt,
  before upper={\setlength{\parskip}{0.5em}}}

\newtcolorbox{algoritmo}[1]{enhanced,breakable,colback=inkc!3,colframe=inkc!60,
  coltitle=white,fonttitle=\bfseries\small,title={#1},
  left=9pt,right=7pt,top=5pt,bottom=5pt,boxrule=0.6pt,arc=2.5pt}

% ---------------------------------------------------------------------------
%  Atajos de notación
% ---------------------------------------------------------------------------
\newcommand{\Zq}{\mathbb{Z}_q}
\newcommand{\mh}{\hat{m}}
\newcommand{\azar}{\xleftarrow{\$}}
\newcommand{\si}{\textcolor{mintc}{\checkmark}}
\newcommand{\no}{\textcolor{coralc}{$\times$}}

% Viñetas explícitas: con babel-spanish y lmodern la de omisión sale como un glifo raro
\setlist[itemize]{label=\textcolor{violetc}{$\bullet$}}

% Opciones del quiz: (a), (b), (c), (d)
\newlist{opciones}{enumerate}{1}
\setlist[opciones]{label=(\alph*),leftmargin=2.4em,itemsep=0.1em,topsep=0.2em}

% ===========================================================================
\begin{document}
% ===========================================================================

% babel-spanish redefine \% inspeccionando \lastskip, lo que revienta dentro de
% los nodos de TikZ. Restauramos la definición estándar de LaTeX.
\makeatletter
\DeclareRobustCommand{\%}{\char`\%\relax}
\makeatother

\begin{titlepage}
\centering
\vspace*{2.5cm}

{\Large\color{violetdeep}\textsc{ChainNotes}\par}
{\small\color{inkc}Módulo 08 · Privacidad\par}

\vspace{1.4cm}
{\Huge\bfseries El anillo de firmas\par}
\vspace{0.4cm}
{\LARGE firmar sin decir quién\par}

\vspace{1.2cm}
\begin{minipage}{0.78\textwidth}
\centering\itshape\color{inkc}
Una firma normal dice «esto lo escribió Ana». Una firma de anillo dice «esto lo
escribió alguien de este grupo», y nadie, ni siquiera con la computadora más
potente del mundo, puede averiguar quién. Aquí lo construimos paso a paso, con
números pequeños y sin necesidad de saber criptografía.
\end{minipage}

\vspace{1.6cm}
\begin{tikzpicture}
  \draw[violetc!60,line width=1pt] (0,0) circle (1.6);
  \foreach \a/\n in {90/A,18/B,-54/C,-126/D,162/E}{
    \fill[white] (\a:1.6) circle (0.42);
    \draw[violetdeep,line width=0.8pt,fill=violetc!15] (\a:1.6) circle (0.38);
    \node[violetdeep,font=\bfseries] at (\a:1.6) {?};
  }
  \node[violetdeep,font=\small] at (0,0) {$\sigma$};
\end{tikzpicture}

\vfill
{\small Apuntes de Blockchain · Proyecto académico\par}
{\small\color{inkc}\today\par}
\end{titlepage}

\tableofcontents
\clearpage

\begin{nota}{Cómo usar este documento}
Es la versión imprimible del apunte interactivo \texttt{temas/anillo-de-firmas.html}.
Cada «laboratorio» de la página web aparece aquí como un ejemplo resuelto con
números fijos, para que puedas seguir todas las cuentas a mano. La sección~\ref{sec:formulas}
reúne la notación formal y es opcional. Al final hay una autoevaluación de 12
preguntas con sus soluciones en el apéndice~\ref{app:soluciones}.
\end{nota}

% ===========================================================================
\section{La idea: firmar en nombre de un grupo}\label{sec:idea}
% ===========================================================================

\begin{nota}{Objetivo}
Entender qué promete una firma de anillo (y qué no) antes de ver una sola fórmula.
\end{nota}

Imagina el tablero de avisos de una facultad. Aparece una carta: «El examen final
tiene un error en la pregunta 4», firmada \textbf{«uno de los cinco profesores de
criptografía»}, con un sello que solo esos cinco pueden poner. Sabes que la carta
es auténtica, porque nadie más tiene ese sello. Pero no sabes cuál de los cinco la
escribió. Eso es, en esencia, una \textbf{firma de anillo}.

Una firma digital ordinaria responde a la pregunta «¿quién escribió esto?» con un
nombre. Una firma de anillo responde «alguien de esta lista». Lo sorprendente son
tres detalles:

\begin{itemize}[leftmargin=1.6em,itemsep=0.3em]
  \item \textbf{El firmante elige el grupo por su cuenta.} Toma las claves públicas
    de otras personas, que cualquiera puede ver, y las mete en su anillo.
  \item \textbf{Los demás ni se enteran.} No firman, no dan permiso y no hacen nada.
    Basta con que su clave pública exista.
  \item \textbf{No hay jefe.} No existe un administrador que pueda «destapar» al
    firmante después. El secreto es para siempre.
\end{itemize}

\begin{figure}[htbp]
\centering
\begin{tikzpicture}[
  font=\small,
  per/.style={circle,draw=#1,fill=#1!15,minimum size=8mm,inner sep=0pt,font=\small\bfseries},
  panel/.style={draw=#1,fill=#1!5,rounded corners=4pt,minimum width=70mm,minimum height=36mm}
]
\node[panel=inkc] (P1) at (0,0) {};
\node[anchor=north,font=\bfseries,inkc] at ($(P1.north)+(0,-2mm)$) {Firma normal};
\foreach \i/\n in {0/A,1/B,2/C,3/D,4/E}{
  \ifnum\i=0
    \node[per=violetdeep] at (-2.4+1.2*\i,0.1) {\textcolor{violetdeep}{\n}};
  \else
    \node[per=inkc,opacity=0.35] at (-2.4+1.2*\i,0.1) {\n};
  \fi
}
\node[font=\small,inkc] at ($(P1.south)+(0,5mm)$) {quien verifica concluye: \textbf{«la firmó Ana»}};

\node[panel=violetc] (P2) at (7.8,0) {};
\node[anchor=north,font=\bfseries,violetdeep] at ($(P2.north)+(0,-2mm)$) {Firma de anillo};
\foreach \i in {0,...,4}{
  \node[per=skyc] at (5.4+1.2*\i,0.1) {\textcolor{skyc!70!black}{?}};
}
\node[font=\small,inkc,align=center] at ($(P2.south)+(0,6mm)$)
  {quien verifica concluye: \textbf{«la firmó alguno}\\\textbf{de Ana, Beto, Carla, Diego o Elena»}};
\end{tikzpicture}
\caption{Lo que aprende quien verifica. En ambos casos sabe que el mensaje es
auténtico; en el anillo, adivinar al firmante al azar acierta 1 vez de cada 5
(20\,\%). Con 10 personas, 1 de cada 10.}
\label{fig:comparacion}
\end{figure}

\begin{clave}{¿Por qué «anillo»?}
No es porque los miembros formen un círculo social. Es por la forma del cálculo: la
firma es una cadena de pasos donde cada uno alimenta al siguiente y el último se
conecta con el primero. Como un círculo no tiene principio, mirar la firma no te
dice dónde empezó y, por lo tanto, quién la hizo.
\end{clave}

\begin{aviso}{Lo que una firma no prueba}
Que la carta esté firmada no significa que lo que dice sea verdad. Solo garantiza
que la escribió alguien del grupo y que nadie la modificó después.
\end{aviso}

El nombre viene de un artículo de 2001 de Rivest, Shamir y Tauman titulado
\emph{How to Leak a Secret} («Cómo filtrar un secreto»): un alto funcionario quiere
denunciar algo y demostrar que es alguien de dentro, sin delatarse.

% ===========================================================================
\section{Las herramientas: claves, reloj y licuadora}\label{sec:claves}
% ===========================================================================

\begin{nota}{Objetivo}
Conocer las tres piezas que usan todos los ejemplos. No hace falta memorizar nada.
\end{nota}

\subsection{Aritmética del reloj (mod)}

En un reloj de 12 horas, 3 horas después de las 11 son las 2, no las 14: al pasar
de 12 vuelves a empezar. Los matemáticos escriben $14 \bmod 12 = 2$ («el residuo de
dividir entre 12»). Aquí usaremos relojes de otros tamaños, como 23 u 11
posiciones. Así los números nunca crecen demasiado.

\subsection{Clave secreta y clave pública}

Cada persona elige un número secreto $x$: su \textbf{clave secreta}. Luego calcula
\[
  y = g^{x} \bmod p
\]
(multiplica $g$ por sí mismo $x$ veces en el reloj de $p$ posiciones) y publica el
resultado: su \textbf{clave pública}.

Hacia adelante la cuenta es fácil. Hacia atrás, dado $y$, encontrar $x$ es muy
difícil: con números de 77 dígitos no hay computadora que lo logre. Es como mezclar
pintura: juntar azul y amarillo es fácil, separarlos después no.

\begin{caso}{Ejemplo con números}
Usamos un reloj de $p = 23$ posiciones y $g = 4$. Ana elige su clave secreta $x = 7$:
\[
  y = 4^{7} \bmod 23 = 16\,384 \bmod 23 = 8 .
\]
Ana publica $y = 8$ y guarda $x = 7$. Las potencias de 4 en este reloj se repiten
cada $q = 11$ pasos:
\begin{center}\small
\begin{tabular}{@{}c*{12}{c}@{}}
\toprule
$k$ & 0 & 1 & 2 & 3 & 4 & 5 & 6 & 7 & 8 & 9 & 10 & 11 \\
\midrule
$4^{k} \bmod 23$ & 1 & 4 & 16 & 18 & 3 & 12 & 2 & 8 & 9 & 13 & 6 & 1 \\
\bottomrule
\end{tabular}
\end{center}
Por eso todo lo que sean exponentes vive en un segundo reloj, de 11 posiciones.
Esta tabla te ahorra casi todas las cuentas de este documento.
\end{caso}

\begin{center}\small
\captionsetup{type=table}
\caption{Los dos relojes y el punto de partida.}
\label{tab:relojes}
\begin{tabular}{@{}l l@{}}
\toprule
\textbf{Símbolo} & \textbf{Papel} \\
\midrule
$p$ (el reloj grande) & donde viven las claves públicas $y$ y los valores $R$ \\
$q$ (el reloj pequeño) & donde viven las claves secretas $x$, los desafíos $c$ y las respuestas $r$ \\
$g$ (el punto de partida) & un número fijo que todos comparten \\
\bottomrule
\end{tabular}
\end{center}

\subsection{La licuadora: una función hash}

Una función hash (módulo 01) mezcla datos y entrega un número que parece
aleatorio: lo mismo entra, lo mismo sale, pero no puedes predecir ni elegir el
resultado. En las firmas, el hash fabrica los \textbf{desafíos}: números que nadie
puede escoger a su conveniencia.

\begin{caso}{Nuestra versión de juguete}
Para que puedas seguir las cuentas a mano, primero el mensaje se resume en un
número pequeño $\mh$, su «huella»: las letras se convierten en un número grande y
nos quedamos con su residuo mód $q$ (la receta exacta está en la
sección~\ref{sec:formulas}). Luego usamos
\[
  H(m, R) = (3\mh + R) \bmod q .
\]
Es una licuadora pésima, pero sirve para ver el mecanismo. Por ejemplo, el mensaje
«hola» tiene huella $\mh = 0$ y «Reunión el viernes» tiene $\mh = 3$ (con $q = 11$).
\end{caso}

\begin{aviso}{Solo para aprender}
Con $q = 11$ cualquiera encuentra $x$ probando 11 opciones. En la vida real $q$
tiene unos 256 bits (unos 77 dígitos).
\end{aviso}

% ===========================================================================
\section{La pieza clave: firmar con la clave o simular sin ella}\label{sec:simular}
% ===========================================================================

\begin{nota}{Objetivo}
Descubrir el truco que hace posible el anillo: se puede «fingir» la firma de
alguien si te dejan escoger las cosas en otro orden.
\end{nota}

Cada eslabón del anillo es un trío de números $(R, c, r)$ que debe cumplir
\textbf{dos condiciones}:
\begin{equation}\label{eq:condiciones}
  \text{\ding{172}}\quad R = g^{r}\cdot y^{c} \bmod p
  \qquad\qquad
  \text{\ding{173}}\quad c = H(m, R).
\end{equation}
La condición \ding{172} amarra el eslabón a la clave pública $y$; la \ding{173} lo
amarra al mensaje. Hay dos formas de conseguir un trío que cumpla \ding{172}:

\begin{description}[leftmargin=1.4em,style=unboxed,itemsep=0.3em]
  \item[Con la clave secreta (orden natural).] Eliges un número al azar $a$,
    calculas $R = g^{a}$, dejas que el hash te dé $c$ y al final respondes
    $r = a - x\cdot c$. Necesitas conocer $x$.
  \item[Sin la clave (simulación, orden invertido).] Escoges primero $c$ y $r$ a tu
    gusto y \emph{despejas} $R = g^{r}\cdot y^{c}$. La condición \ding{172} se cumple
    siempre, sin saber $x$.
\end{description}

\begin{figure}[htbp]
\centering
\begin{tikzpicture}[
  font=\small,
  paso/.style={draw=#1,fill=#1!10,rounded corners=3pt,minimum height=9mm,
               minimum width=27mm,align=center,inner sep=3pt},
  >={Stealth[length=2.2mm]}
]
\node[font=\bfseries,mintc,anchor=west] at (-1.6,1.1) {Con clave};
\node[paso=mintc] (a1) at (0,0) {$a$ al azar};
\node[paso=mintc,right=7mm of a1] (a2) {$R = g^{a}$};
\node[paso=mintc,right=7mm of a2] (a3) {$c = H(m,R)$};
\node[paso=mintc,right=7mm of a3] (a4) {$r = a - x c$\\[-1pt]{\scriptsize usa $x$}};
\draw[->,mintc,thick] (a1) -- (a2);
\draw[->,mintc,thick] (a2) -- (a3);
\draw[->,mintc,thick] (a3) -- (a4);

\node[font=\bfseries,amberc,anchor=west] at (-1.6,-1.6) {Simulación};
\node[paso=amberc] (b1) at (0,-2.7) {eliges $c$ y $r$};
\node[paso=amberc,right=7mm of b1] (b2) {$R = g^{r} y^{c}$};
\node[paso=amberc,right=7mm of b2] (b3) {$H(m,R) \overset{?}{=} c$};
\node[paso=coralc,right=7mm of b3] (b4) {casi nunca\\[-1pt]{\scriptsize 1 de cada $q$}};
\draw[->,amberc,thick] (b1) -- (b2);
\draw[->,amberc,thick] (b2) -- (b3);
\draw[->,coralc,thick] (b3) -- (b4);
\end{tikzpicture}
\caption{Las dos maneras de llenar un eslabón. Arriba, el hash se evalúa
\emph{antes} de la respuesta y la clave hace encajar todo. Abajo, el desafío se
escogió \emph{antes} que $R$, así que el hash de ese $R$ casi nunca coincide.}
\label{fig:dosordenes}
\end{figure}

¿Entonces cualquiera puede falsificar? No, porque la simulación casi nunca cumple
la condición \ding{173}: tú escogiste $c$ antes de conocer $R$, y el hash de ese $R$
sale otro número. Es como contestar un examen inventando tú las preguntas: todas
tus respuestas encajan, pero no con las preguntas oficiales.

\begin{proposicion}[Por qué funciona con la clave]\label{prop:schnorr}
Si $y = g^{x}$ y $r = a - x c \bmod q$, entonces $g^{r} y^{c} = g^{a} = R$.
\end{proposicion}
\begin{proof}
$g^{r} y^{c} = g^{a - xc}\,(g^{x})^{c} = g^{a - xc}\, g^{xc} = g^{a}$. El término
$xc$ aparece restando y sumando, y se cancela. (Los exponentes se pueden reducir
módulo $q$ porque $g^{q} = 1$.)
\end{proof}

Esta es la firma de \textbf{Schnorr} (1989), una de las más usadas hoy: Bitcoin la
adoptó en 2021 con la actualización Taproot.

\begin{laboratorio}{Laboratorio 1 — Firma real contra simulación}
Clave de Ana: $x = 7$, $y = 8$. Grupo $p = 23$, $q = 11$, $g = 4$. Mensaje «hola»,
con huella $\mh = 0$.

\textbf{Con la clave}, eligiendo $a = 6$:
\begin{enumerate}[leftmargin=2em,itemsep=0.1em,topsep=0.2em]
  \item $R = 4^{6} \bmod 23 = 2$.
  \item $c = H(m, R) = (3\cdot 0 + 2) \bmod 11 = 2$.
  \item $r = a - x c = 6 - 7\cdot 2 = -8 \equiv 3 \pmod{11}$.
\end{enumerate}
Comprobación: \ding{172} $4^{3}\cdot 8^{2} \bmod 23 = 18\cdot 18 \bmod 23 = 2 = R$
\si; \ding{173} $c = H(m,R)$ \si. Firma válida $(c, r) = (2, 3)$.

\textbf{Simulando}, escogiendo $c = 6$ y $r = 5$:
\begin{enumerate}[leftmargin=2em,itemsep=0.1em,topsep=0.2em]
  \item $R = 4^{5}\cdot 8^{6} \bmod 23 = 12 \cdot 13 \bmod 23 = 18$.
  \item El hash de ese $R$: $H(m, 18) = 18 \bmod 11 = 7$.
\end{enumerate}
Comprobación: \ding{172} se cumple por construcción \si; \ding{173} $6 \neq 7$ \no.
No es una firma.

\textbf{Qué observar.} Con $r = 5$, de los 11 desafíos posibles solo $c = 3$ hace
coincidir el hash. Para cada $r$ suele haber uno o dos valores que funcionan (a
veces ninguno), y solo se encuentran \emph{probando}; con un $q$ de 77 dígitos eso es
imposible. En la página web puedes mover $a$, $c$ y $r$ y ver las cuentas en vivo.
\end{laboratorio}

\begin{clave}{La idea que hay que llevarse}
La simulación funciona si el desafío $c$ te llega \textbf{de fuera}, ya fijado, y tú
solo pones la respuesta. En el anillo pasa justo eso: cada desafío lo produce el
eslabón anterior. Así el firmante puede rellenar los eslabones de los demás sin
conocer sus claves.
\end{clave}

% ===========================================================================
\section{Construye un anillo: una cadena que se muerde la cola}\label{sec:anillo}
% ===========================================================================

\begin{nota}{Objetivo}
Ver cómo el firmante rellena los eslabones de los demás y cierra el último con su
propia clave.
\end{nota}

Piensa en una carrera de relevos en una pista circular. Cada corredor recibe el
testigo (el desafío $c$) del anterior, hace su parte y produce el testigo del
siguiente. El firmante organiza toda la carrera él solo, en tres fases:

\begin{enumerate}[leftmargin=2em,itemsep=0.4em,label=\textbf{\arabic*.}]
  \item \textbf{Arranca en su propio lugar.} Elige un número al azar $a$, calcula su
    $R = g^{a}$ y el hash le da el desafío para la persona siguiente.
  \item \textbf{Simula a todos los demás.} A cada miembro le llega su desafío ya
    hecho. El firmante escoge una respuesta $r$ al azar, despeja su $R$ y obtiene el
    desafío del siguiente. No necesita sus claves secretas.
  \item \textbf{Cierra el hueco con su clave.} La vuelta regresa a él con un desafío
    que no escogió. Para que encaje con el $R$ que fijó al principio calcula
    $r = a - x\cdot c$. Es el único paso que exige una clave secreta.
\end{enumerate}

La firma que se publica es muy compacta: el primer desafío $c_1$ y una respuesta
$r$ por cada miembro.
\begin{equation}\label{eq:sigma}
  \sigma = (c_1;\ r_1, r_2, \dots, r_n)
  \qquad\text{($n+1$ números para un anillo de $n$ personas).}
\end{equation}
Los $R$ no se publican, porque cualquiera puede recalcularlos. Junto con el mensaje
y la lista de claves públicas del anillo, eso es todo lo que necesita quien verifica.

\begin{laboratorio}{Laboratorio 2 — Un anillo de cinco, paso a paso}
Grupo $p = 23$, $q = 11$, $g = 4$. Mensaje «Reunión el viernes», huella $\mh = 3$,
así que $H(m,R) = (9 + R) \bmod 11$. Firma \textbf{Carla} (miembro 3).

\begin{center}\small
\begin{tabular}{@{}l c c l@{}}
\toprule
\textbf{Miembro} & \textbf{$x$ (secreta)} & \textbf{$y$ (pública)} & \\
\midrule
1 · Ana   & \textcolor{inkc!45}{7} & 8  & \\
2 · Beto  & \textcolor{inkc!45}{5} & 12 & \\
3 · Carla & 3 & 18 & \textcolor{violetdeep}{firma} \\
4 · Diego & \textcolor{inkc!45}{9} & 13 & \\
5 · Elena & \textcolor{inkc!45}{4} & 3  & \\
\bottomrule
\end{tabular}
\end{center}
Solo Carla conoce su $x$. Las de los demás van en gris porque en la vida real nadie
las ve; solo las usaremos en la sección~\ref{sec:quien}.

\begin{description}[leftmargin=1.2em,style=unboxed,itemsep=0.35em,font=\normalfont\bfseries\color{violetdeep}]
  \item[Paso 1 · Carla arranca.] Elige al azar $a = 6$ y calcula
    $R_3 = 4^{6} \bmod 23 = 2$. El hash da el desafío de Diego:
    $c_4 = H(m, 2) = 11 \bmod 11 = 0$.
  \item[Paso 2 · Diego se simula.] Le llega $c_4 = 0$; se escoge $r_4 = 5$ y se
    despeja $R_4 = 4^{5}\cdot 13^{0} \bmod 23 = 12$. Desafío de Elena:
    $c_5 = H(m,12) = 21 \bmod 11 = 10$.
  \item[Paso 3 · Elena se simula.] Le llega $c_5 = 10$; se escoge $r_5 = 8$ y
    $R_5 = 4^{8}\cdot 3^{10} \bmod 23 = 9\cdot 8 \bmod 23 = 3$. Desafío de Ana:
    $c_1 = H(m,3) = 12 \bmod 11 = 1$.
  \item[Paso 4 · Ana se simula.] Le llega $c_1 = 1$; se escoge $r_1 = 2$ y
    $R_1 = 4^{2}\cdot 8^{1} \bmod 23 = 128 \bmod 23 = 13$. Desafío de Beto:
    $c_2 = H(m,13) = 22 \bmod 11 = 0$.
  \item[Paso 5 · Beto se simula.] Le llega $c_2 = 0$; se escoge $r_2 = 10$ y
    $R_2 = 4^{10}\cdot 12^{0} \bmod 23 = 6$. Desafío de Carla:
    $c_3 = H(m,6) = 15 \bmod 11 = 4$.
  \item[Paso 6 · Carla cierra.] La vuelta le regresa $c_3 = 4$, un desafío que no
    escogió. Con su clave $x = 3$:
    $r_3 = a - x c_3 = 6 - 3\cdot 4 = -6 \equiv 5 \pmod{11}$. \textcolor{mintc}{\textbf{Anillo cerrado.}}
\end{description}

La firma publicada es
\[
  \sigma = (c_1 = 1;\ r_1 = 2,\ r_2 = 10,\ r_3 = 5,\ r_4 = 5,\ r_5 = 8).
\]
Tamaño: 6 números. Claves públicas que necesita el verificador: 5. Probabilidad de
adivinar al firmante: $1/5$.
\end{laboratorio}

\begin{figure}[htbp]
\centering
\begin{tikzpicture}[
  font=\small,
  nodo/.style={circle,draw=#1!80!black,fill=#1,minimum size=15mm,inner sep=0pt,
               line width=0.9pt,align=center},
  >={Stealth[length=2.4mm]}
]
\def\RR{3.1}
\draw[inkc!10,line width=10pt] (0,0) circle (\RR);
% arcos en sentido horario, del miembro i al i+1, con el desafío que llevan
% (nombres \angA… para no pisar \aa, \t o \c, que en LaTeX son acentos)
\foreach \i/\cl in {0/{c_2 = 0},1/{c_3 = 4},2/{c_4 = 0},3/{c_5 = 10},4/{c_1 = 1}}{
  \pgfmathsetmacro{\angA}{90-72*\i-15}
  \pgfmathsetmacro{\angB}{90-72*(\i+1)+15}
  \pgfmathsetmacro{\angM}{90-72*\i-36}
  \draw[->,violetdeep,line width=1.1pt] (\angA:\RR) arc[start angle=\angA,end angle=\angB,radius=\RR];
  \node[violetdeep,font=\small\bfseries,fill=white,inner sep=1.5pt] at (\angM:\RR-0.85) {$\cl$};
}
% nodos
\foreach \i/\n/\y/\col in {0/Ana/8/mintc,1/Beto/12/mintc,2/Carla/18/violetc,3/Diego/13/mintc,4/Elena/3/mintc}{
  \pgfmathsetmacro{\angN}{90-72*\i}
  \node[nodo=\col] at (\angN:\RR) {\textbf{\textcolor{white}{\n}}\\[-2pt]{\scriptsize\textcolor{white}{$y = \y$}}};
}
% valores R y r fuera del anillo
\node[inkc,align=center,font=\footnotesize] at (90:\RR+1.25) {$R_1 = 13$, $r_1 = 2$};
\node[inkc,align=left,anchor=west,font=\footnotesize] at (18:\RR+0.9) {$R_2 = 6$\\$r_2 = 10$};
\node[violetdeep,align=left,anchor=west,font=\footnotesize] at (-54:\RR+0.9) {$a = 6$, $R_3 = 2$\\$r_3 = 5$ \textbf{(cierre)}};
\node[inkc,align=right,anchor=east,font=\footnotesize] at (-126:\RR+0.9) {$R_4 = 12$\\$r_4 = 5$};
\node[inkc,align=right,anchor=east,font=\footnotesize] at (162:\RR+0.9) {$R_5 = 3$\\$r_5 = 8$};
\node[inkc!70,font=\footnotesize,align=center] at (0,0) {el cálculo\\da la vuelta\\en sentido horario};
\end{tikzpicture}
\caption{El anillo del laboratorio 2 ya cerrado. Violeta: la firmante (Carla), que
arrancó en su lugar con $a = 6$ y cerró con su clave. Verde: los miembros
simulados. Cada flecha lleva el desafío que un eslabón le entrega al siguiente.}
\label{fig:anillo}
\end{figure}

\begin{aviso}{El precio del anonimato}
La firma crece con el tamaño del anillo: esconderte entre 1000 personas cuesta una
firma 1000 veces más larga. Existen variantes más compactas, pero esta es la idea
básica.
\end{aviso}

% ===========================================================================
\section{Verificar: dar la vuelta completa}\label{sec:verificar}
% ===========================================================================

\begin{nota}{Objetivo}
Comprobar una firma tú mismo y ver que cualquier cambio, por pequeño que sea, rompe
el anillo.
\end{nota}

Quien verifica no conoce ninguna clave secreta ni sabe quién firmó. Tiene el
mensaje, las claves públicas del anillo y la firma $\sigma$. Su trabajo es
sencillo:

\begin{enumerate}[leftmargin=2em,itemsep=0.2em]
  \item Empieza con el desafío $c_1$ que viene en la firma.
  \item Para cada miembro, recalcula su $R$ con la condición \ding{172} y, con el
    hash, el desafío del siguiente.
  \item Al terminar la vuelta compara: si vuelve a obtener exactamente $c_1$, la
    firma es \textbf{válida}.
\end{enumerate}
En símbolos:
\begin{equation}\label{eq:verificar}
  \text{para } i = 1,\dots,n:\quad R_i = g^{r_i} y_i^{c_i},\quad c_{i+1} = H(m, R_i);
  \qquad \text{válida} \iff c_{n+1} = c_1 .
\end{equation}

Es como seguir instrucciones del tipo «da 3 pasos, gira, da 5 pasos…» en un
círculo: si al final llegas al punto de partida, las instrucciones son correctas.
Si alguien cambió una sola, terminas en otro lugar.

Fíjate en que la regla trata a todos los miembros exactamente igual: no hay ningún
paso especial para el firmante. Esa simetría es la raíz del anonimato.

\begin{laboratorio}{Laboratorio 3 — Verifica y rompe}
Verificamos la firma del laboratorio 2, $\sigma = (1;\ 2, 10, 5, 5, 8)$, con el
mensaje original ($\mh = 3$) y con uno alterado: «Reunión el \emph{jueves}»
($\mh = 6$, así que $H(m,R) = (18 + R) \bmod 11$).

\begin{center}\small
\begin{tabular}{@{}l c c c c c@{}}
\toprule
& \multicolumn{2}{c}{\textbf{Mensaje original}} & & \multicolumn{2}{c}{\textbf{Mensaje alterado}} \\
\cmidrule(lr){2-3}\cmidrule(lr){5-6}
\textbf{Eslabón} & $R_i$ & $c_{i+1}$ & & $R_i$ & $c_{i+1}$ \\
\midrule
Ana   ($c_1 = 1$)   & $4^{2}\cdot 8^{1} = 13$    & 0  & & $13$ & 9 \\
Beto  & $4^{10}\cdot 12^{0} = 6$  & 4  & & $4^{10}\cdot 12^{9} = 1$  & 8 \\
Carla & $4^{5}\cdot 18^{4} = 2$   & 0  & & $4^{5}\cdot 18^{8} = 8$   & 4 \\
Diego & $4^{5}\cdot 13^{0} = 12$  & 10 & & $4^{5}\cdot 13^{4} = 9$   & 5 \\
Elena & $4^{8}\cdot 3^{10} = 3$   & \textbf{1} & & $4^{8}\cdot 3^{5} = 2$ & \textbf{9} \\
\midrule
\textbf{Resultado} & \multicolumn{2}{c}{$c_6 = 1 = c_1$ \si\ válida} & & \multicolumn{2}{c}{$c_6 = 9 \neq 1$ \no\ inválida} \\
\bottomrule
\end{tabular}
\end{center}

Observa que en el eslabón de Carla, la firmante, el verificador hace la misma
cuenta que en los demás: $4^{5}\cdot 18^{4} \bmod 23 = 12\cdot 4 \bmod 23 = 2$.
Nada lo distingue.

\textbf{Qué observar.} Cambiar una letra cambia la huella, que cambia todos los
desafíos. Cambiar una respuesta $r$ cambia un $R$, que cambia el desafío siguiente,
y el error se arrastra por el resto de la vuelta. En la web puedes editar cualquier
número y ver el resultado en vivo.
\end{laboratorio}

\begin{aviso}{Coincidencias por suerte}
Con un reloj de solo $q = 11$ posiciones, uno o dos de cada $q$ cambios hacen que la
vuelta cierre por casualidad. Por eso en la práctica $q$ tiene 77 dígitos y la
suerte deja de existir.
\end{aviso}

% ===========================================================================
\section{¿Quién firmó? Cualquiera pudo hacerlo}\label{sec:quien}
% ===========================================================================

\begin{nota}{Objetivo}
Convencerte de que la firma no delata a nadie, ni siquiera ante alguien con poder de
cómputo ilimitado.
\end{nota}

Hagamos de detectives. Tomamos la firma válida y nos preguntamos: «¿pudo haberla
hecho Ana?». Si Ana hubiera firmado, habría empezado en su propio eslabón con un
número al azar $a'$ que cumpliera $g^{a'} = R_{\text{Ana}}$, y habría simulado a
todos los demás con las mismas respuestas que vemos en la firma. ¿Existe ese $a'$?
Sí:
\begin{equation}\label{eq:detective}
  a' = r_i + x_i c_i \bmod q
  \qquad\Longrightarrow\qquad
  g^{a'} = g^{r_i}\cdot (g^{x_i})^{c_i} = g^{r_i} y_i^{c_i} = R_i .
\end{equation}

Lo mismo vale para Beto, para Carla y para todos. \textbf{Para cada miembro existe un
camino válido que produce exactamente la misma firma}, y todos esos caminos son
igual de probables: cada $r$ de los simulados se escogió al azar, igual que el $a$
del firmante. La firma no favorece a nadie.

\begin{laboratorio}{Laboratorio 4 — Todos son sospechosos}
Con la firma del laboratorio 2, suponemos uno por uno que firmó cada miembro y
reconstruimos su número inicial. (Esta tabla usa las claves secretas de todos, algo
que nadie tiene en la práctica.)

\begin{center}\small
\begin{tabular}{@{}l c c c l c c c@{}}
\toprule
\textbf{Si firmó…} & $x$ & $c$ recibido & $r$ & $a' = r + x c \bmod 11$ & $g^{a'}$ & $R$ del anillo & ¿posible? \\
\midrule
Ana   & 7 & 1  & 2  & $2 + 7\cdot 1 = 9$            & 13 & 13 & \si \\
Beto  & 5 & 0  & 10 & $10 + 5\cdot 0 = 10$          & 6  & 6  & \si \\
\textcolor{violetdeep}{\textbf{Carla}} & 3 & 4  & 5  & $5 + 3\cdot 4 = 17 \equiv 6$  & 2  & 2  & \si\ \textcolor{violetdeep}{\textbf{(firmó)}} \\
Diego & 9 & 0  & 5  & $5 + 9\cdot 0 = 5$            & 12 & 12 & \si \\
Elena & 4 & 10 & 8  & $8 + 4\cdot 10 = 48 \equiv 4$ & 3  & 3  & \si \\
\bottomrule
\end{tabular}
\end{center}

Para Carla sale $a' = 6$: justo el número que de verdad usó. Pero las otras cuatro
filas funcionan igual de bien. Ni con toda la información del mundo se distingue a la
firmante.
\end{laboratorio}

\begin{clave}{Anonimato incondicional}
No es que averiguar al firmante sea «muy difícil»: es que la información no está en
la firma. Ni con una computadora infinitamente rápida, ni conociendo todas las
claves secretas, se puede saber quién fue. Lo mejor que puede hacer un detective es
adivinar: acierta 1 vez de cada $n$.
\end{clave}

% ===========================================================================
\section{Para qué sirve: Monero y otras aplicaciones}\label{sec:blockchain}
% ===========================================================================

\begin{nota}{Objetivo}
Ver cómo una blockchain usa las firmas de anillo para ocultar quién paga sin abrir
la puerta al doble gasto.
\end{nota}

En Bitcoin todo es público: cada pago muestra qué monedas se gastan y quién las
firma. Cualquiera puede seguir el rastro del dinero. \textbf{Monero}, una
criptomoneda lanzada en 2014, usa firmas de anillo para romper ese rastro.

\subsection{Los señuelos}

Cuando gastas una moneda en Monero, tu billetera toma otras monedas que ya existen
en la cadena, de otras personas, y las mete en el anillo como \textbf{señuelos}.
Desde 2022 cada anillo tiene 16 miembros: tu moneda real y 15 señuelos. Quien mira
la cadena sabe que se gastó una de esas 16, pero no cuál.

\subsection{El problema: ¿y si gasto dos veces?}

Si nadie sabe qué moneda se gastó, ¿cómo evita la red el doble gasto (módulo 07)?
Con una variante llamada \textbf{firma de anillo enlazable}. Cada firma incluye una
\textbf{imagen de clave}: una especie de «huella» que se calcula con la clave
secreta de la moneda. No revela a cuál miembro del anillo pertenece, pero es
\emph{siempre la misma} para la misma moneda. Si alguien intenta gastarla otra vez,
aparece la misma huella y la red rechaza la transacción.

\begin{figure}[htbp]
\centering
\begin{tikzpicture}[
  font=\small,
  moneda/.style={circle,draw=skyc!80!black,fill=skyc!15,minimum size=6.5mm,inner sep=0pt,font=\scriptsize},
  caja/.style={draw=#1,fill=#1!8,rounded corners=3pt,align=center,inner sep=5pt},
  >={Stealth[length=2.2mm]}
]
\foreach \k in {0,...,15}{
  \node[moneda] (m\k) at ({90-22.5*\k}:2.0) {?};
}
\node[align=center,font=\footnotesize,inkc] at (0,0) {16 monedas:\\1 real + 15\\señuelos};
\node[caja=violetc] (I) at (6.2,1.0) {\textbf{imagen de clave} $I$\\{\scriptsize igual cada vez que se gasta la misma moneda}};
\node[caja=inkc] (L) at (6.2,-1.0) {lista de imágenes ya usadas\\{\scriptsize ¿ya apareció $I$?}};
\draw[->,violetdeep,thick] (2.5,0.6) -- (I.west);
\draw[->,inkc,thick] (I) -- (L);
\node[caja=mintc,font=\scriptsize] at (4.2,-2.6) {no: el pago es válido};
\node[caja=coralc,font=\scriptsize] at (8.2,-2.6) {sí: doble gasto, se rechaza};
\draw[->,mintc,thick] (L.south) -- ++(-1.2,-0.75);
\draw[->,coralc,thick] (L.south) -- ++(1.2,-0.75);
\end{tikzpicture}
\caption{El anillo oculta cuál moneda se gasta; la imagen de clave delata un
segundo gasto de la misma moneda sin revelar cuál de las 16 era.}
\label{fig:monero}
\end{figure}

\begin{caso}{Tres capas de privacidad en Monero}
Las firmas de anillo ocultan \textbf{quién paga}. Las direcciones ocultas
(\emph{stealth addresses}) ocultan \textbf{quién recibe}. Y las transacciones
confidenciales (RingCT) ocultan \textbf{cuánto se paga} usando compromisos de
Pedersen (módulo 02).
\end{caso}

\subsection{Otras aplicaciones}

\begin{itemize}[leftmargin=1.6em,itemsep=0.3em]
  \item \textbf{Denuncias anónimas.} Un empleado demuestra que trabaja en la empresa
    sin decir quién es: el caso que dio origen a la idea.
  \item \textbf{Votaciones.} Con firmas enlazables, cada votante prueba que está en el
    padrón y la red detecta si alguien vota dos veces, sin saber quién votó qué.
  \item \textbf{Pertenencia a un grupo.} Probar «soy miembro» (de un club, de una
    universidad) sin revelar cuál miembro.
\end{itemize}

\begin{aviso}{Los límites del anonimato}
Las matemáticas son perfectas, pero el anonimato depende de los señuelos. Si un
señuelo es obviamente falso (por ejemplo, una moneda que se sabe gastada en otro
lado), se descarta y el anillo «encoge». En los primeros años de Monero se permitían
anillos sin señuelos, y estudios de 2017 y 2018 lograron rastrear muchas de esas
transacciones. Por eso hoy el tamaño es obligatorio y los señuelos se eligen con más
cuidado.
\end{aviso}

\begin{center}\small
\captionsetup{type=table}
\caption{Cuatro tipos de firma frente a frente.}
\label{tab:tipos}
\begin{tabularx}{\textwidth}{@{}>{\raggedright\arraybackslash}p{2.7cm} >{\raggedright\arraybackslash}X >{\raggedright\arraybackslash}X >{\raggedright\arraybackslash}X >{\raggedright\arraybackslash}X@{}}
\toprule
\textbf{Tipo} & \textbf{¿Revela quién firmó?} & \textbf{¿Hay administrador?} & \textbf{¿Los demás participan?} & \textbf{Tamaño} \\
\midrule
Normal & \textcolor{coralc}{Sí} & No & No hay «demás» & Fijo y pequeño \\
Multifirma (varias personas firman juntas) & \textcolor{amberc}{Normalmente sí} & No & Sí, todas deben firmar & Crece con quienes firman (o fijo, en versiones agregadas) \\
De grupo & \textcolor{amberc}{Solo al administrador} & Sí: inscribe miembros y puede destapar al firmante & Deben inscribirse antes & Fijo \\
De anillo & \textcolor{mintc}{A nadie, nunca} & No & No, ni se enteran & Crece con el anillo \\
\bottomrule
\end{tabularx}
\end{center}

\begin{center}\small
\captionsetup{type=table}
\caption{Línea de tiempo.}
\label{tab:tiempo}
\begin{tabular}{@{}r l@{}}
\toprule
\textbf{Año} & \textbf{Hito} \\
\midrule
1989 & Schnorr publica su esquema de firma. \\
2001 & Rivest, Shamir y Tauman inventan las firmas de anillo (con RSA). \\
2002 & Abe, Ohkubo y Suzuki dan la versión basada en Schnorr, la de este apunte. \\
2004 & Liu, Wei y Wong proponen las firmas de anillo enlazables. \\
2013 & El protocolo CryptoNote las lleva a las criptomonedas. \\
2014 & Nace Monero, basado en CryptoNote. \\
2017 & RingCT oculta también los montos. \\
2020 & Monero adopta CLSAG, firmas de anillo más pequeñas y rápidas. \\
2022 & El anillo pasa a 16 miembros. \\
\bottomrule
\end{tabular}
\end{center}

% ===========================================================================
\section{Rincón técnico: el esquema completo}\label{sec:formulas}
% ===========================================================================

\begin{nota}{Opcional}
Todo lo anterior en notación formal, para quien quiera ir más allá. Puedes saltar a
la síntesis (sección~\ref{sec:sintesis}).
\end{nota}

Esta es la versión de \textbf{Abe, Ohkubo y Suzuki (2002)}, basada en firmas de
Schnorr. La construcción original de Rivest, Shamir y Tauman (2001) usa RSA, pero la
idea del anillo es la misma. Los índices son cíclicos: después de $n$ viene $1$.

\begin{definicion}[Parámetros y claves]
Sea $p$ primo y $g \in \mathbb{Z}_p^{*}$ de orden primo $q$, con $q \mid p-1$. Cada
miembro $i$ tiene una clave secreta $x_i \in \Zq$ y una pública
$y_i = g^{x_i} \bmod p$. El anillo es la lista $L = (y_1,\dots,y_n)$ y
$H : \{0,1\}^{*} \to \Zq$ es una función hash; los desafíos son
$c = H(L, m, R)$. Incluir $L$ en el hash amarra la firma a ese anillo concreto.
\end{definicion}

\begin{algoritmo}{Firmar (miembro $s$)}
\begin{enumerate}[leftmargin=1.6em,itemsep=0.15em,label=\arabic*:]
  \item $a \azar \Zq$
  \item $R_s \gets g^{a} \bmod p$, \quad $c_{s+1} \gets H(L, m, R_s)$
  \item \textbf{para} $i = s+1, \dots, s-1$ (cíclico):
  \begin{itemize}[leftmargin=1.4em,label={}]
    \item $r_i \azar \Zq$, \quad $R_i \gets g^{r_i} y_i^{c_i} \bmod p$, \quad
          $c_{i+1} \gets H(L, m, R_i)$
  \end{itemize}
  \item $r_s \gets a - x_s c_s \bmod q$
  \item \textbf{devolver} $\sigma = (c_1;\ r_1, \dots, r_n)$
\end{enumerate}
\end{algoritmo}

\begin{algoritmo}{Verificar}
\begin{enumerate}[leftmargin=1.6em,itemsep=0.15em,label=\arabic*:]
  \item \textbf{para} $i = 1, \dots, n$: \quad $R_i \gets g^{r_i} y_i^{c_i} \bmod p$,
        \quad $c_{i+1} \gets H(L, m, R_i)$
  \item \textbf{aceptar} si y solo si $c_{n+1} = c_1$
\end{enumerate}
\end{algoritmo}

\begin{proposicion}[Correctitud]\label{prop:correcta}
Una firma producida honestamente siempre verifica.
\end{proposicion}
\begin{proof}
En los nodos $i \neq s$, el verificador recalcula $R_i$ con la misma fórmula con que
se definió, así que obtiene los mismos $c_{i+1}$. En el nodo del firmante, por la
proposición~\ref{prop:schnorr},
\[
  g^{r_s} y_s^{c_s} = g^{a - x_s c_s}\, g^{x_s c_s} = g^{a} = R_s ,
\]
de modo que el hash entrega el mismo $c_{s+1}$ y la vuelta se completa:
$c_{n+1} = c_1$.
\end{proof}

\begin{proposicion}[Anonimato incondicional]\label{prop:anonimato}
Para un mensaje y un anillo fijos, la distribución de $\sigma$ no depende de $s$.
\end{proposicion}
\begin{proof}[Idea de la prueba]
Fijado $c_1$, la firma queda determinada por las $n$ respuestas. Las $n-1$
respuestas simuladas son uniformes en $\Zq$ por construcción, y $r_s = a - x_s c_s$
también lo es porque $a$ es uniforme e independiente. Por la
ecuación~\eqref{eq:detective}, cada $\sigma$ válida se obtiene desde cualquier $s$
con exactamente una elección de azar. Luego todas las $\sigma$ válidas son
equiprobables sea quien sea el firmante.
\end{proof}

\begin{observacion}[No falsificabilidad]
Sin conocer ninguna $x_i$ del anillo, cerrar la vuelta exige acertar un hash o
calcular un logaritmo discreto. En el modelo de oráculo aleatorio se demuestra que
falsificar es tan difícil como resolver el logaritmo discreto en el grupo.
\end{observacion}

\begin{caso}{Diferencias con los ejemplos de este documento}
Aquí usamos $H(m,R) = (3\mh + R) \bmod q$, que ignora $L$, y grupos diminutos
($p = 23$, $q = 11$; la página web ofrece también $p = 47$, $q = 23$ y $p = 59$,
$q = 29$). La huella del mensaje se calcula así: con $u_1,\dots,u_\ell$ los códigos
de sus caracteres (UTF-16),
\[
  h_0 = 7,\qquad h_k = (131\, h_{k-1} + u_k) \bmod 1\,000\,003,\qquad \mh = h_\ell \bmod q .
\]
En la práctica $q$ tiene unos 256 bits (normalmente sobre curvas elípticas) y $H$ es
un hash criptográfico como SHA-256.
\end{caso}

% ===========================================================================
\section{Síntesis}\label{sec:sintesis}
% ===========================================================================

\begin{clave}{Lo que hay que llevarse}
\begin{itemize}[leftmargin=1.4em,itemsep=0.3em]
  \item Una firma de anillo prueba que \textbf{alguien de un grupo} firmó, sin revelar
    quién. El firmante arma el grupo solo, con claves públicas ajenas, y no hay
    administrador.
  \item El truco es la \textbf{simulación}: si el desafío te llega ya hecho, puedes
    escoger la respuesta y despejar el resto sin conocer la clave.
  \item El firmante simula a todos los demás y usa su clave secreta \textbf{una sola
    vez}, para cerrar la vuelta en su propio eslabón.
  \item El verificador recorre el anillo y comprueba que regresa a $c_1$. Cualquier
    cambio rompe el cierre.
  \item El anonimato es \textbf{incondicional}: todo miembro tiene un camino válido
    hacia la misma firma. Su costo es una firma que crece con el anillo.
  \item Monero usa una versión \textbf{enlazable}: la imagen de clave impide el doble
    gasto sin revelar al firmante.
\end{itemize}
\end{clave}

\begin{center}\small
\captionsetup{type=table}
\caption{Glosario mínimo.}
\label{tab:glosario}
\begin{tabularx}{\textwidth}{@{}l X@{}}
\toprule
\textbf{Término} & \textbf{Significado} \\
\midrule
Anillo & Lista de claves públicas entre las que se esconde el firmante. \\
Desafío $c$ & Número producido por el hash; nadie puede escogerlo a su gusto. \\
Respuesta $r$ & Número que publica cada eslabón para que la condición \ding{172} se cumpla. \\
Simulación & Llenar un eslabón sin clave: escoger $r$ y despejar $R$. \\
Señuelo & Miembro del anillo que no firmó; en Monero, una moneda ajena. \\
Imagen de clave & Huella de la clave secreta que delata un segundo gasto, pero no al firmante. \\
\bottomrule
\end{tabularx}
\end{center}

\begin{center}\small
\captionsetup{type=table}
\caption{Todo el módulo en seis renglones.}
\label{tab:resumen}
\begin{tabular}{@{}l l l@{}}
\toprule
\textbf{Concepto} & \textbf{Expresión} & \textbf{Dónde} \\
\midrule
Clave pública & $y = g^{x} \bmod p$ & §\ref{sec:claves} \\
Eslabón & $R = g^{r} y^{c}$, \ $c = H(m,R)$ & Ec.~\eqref{eq:condiciones} \\
Cierre del firmante & $r_s = a - x_s c_s \bmod q$ & §\ref{sec:anillo} \\
Firma publicada & $\sigma = (c_1;\ r_1,\dots,r_n)$ & Ec.~\eqref{eq:sigma} \\
Verificación & válida $\iff c_{n+1} = c_1$ & Ec.~\eqref{eq:verificar} \\
Camino del detective & $a' = r_i + x_i c_i$, \ $g^{a'} = R_i$ & Ec.~\eqref{eq:detective} \\
\bottomrule
\end{tabular}
\end{center}

% ===========================================================================
\section{Pon a prueba lo aprendido}\label{sec:quiz}
% ===========================================================================

\begin{nota}{Instrucciones}
Doce preguntas sobre los conceptos clave, cada una con una sola respuesta correcta.
Contéstalas antes de mirar las soluciones del apéndice~\ref{app:soluciones}, que
explican cada respuesta e indican qué sección repasar. En la página web el quiz se
corrige solo.
\end{nota}

\begin{enumerate}[leftmargin=2em,itemsep=0.9em,label=\textbf{\arabic*.}]
\item Una firma de anillo válida sobre un mensaje le demuestra a quien la verifica que…
  \begin{opciones}
    \item lo firmó la persona que aparece primero en el anillo.
    \item lo firmó alguno de los miembros del anillo, sin revelar cuál.
    \item lo firmaron todos los miembros del anillo juntos.
    \item lo que dice el mensaje es verdad.
  \end{opciones}

\item ¿Qué necesita la persona que firma para armar el anillo?
  \begin{opciones}
    \item Las claves secretas de todos los miembros.
    \item El permiso de cada miembro y la aprobación de un administrador.
    \item Su propia clave secreta y las claves públicas de los demás.
    \item Solo claves públicas; nadie usa claves secretas.
  \end{opciones}

\item ¿En qué se diferencia una firma de anillo de una firma de grupo?
  \begin{opciones}
    \item En la de grupo hay un administrador que puede revelar quién firmó; en la de
      anillo nadie tiene ese poder.
    \item La firma de anillo siempre es más pequeña.
    \item En la de anillo los miembros deben inscribirse antes.
    \item No hay diferencia: son dos nombres para lo mismo.
  \end{opciones}

\item Al simular el eslabón de otra persona, ¿cuál es el truco?
  \begin{opciones}
    \item Adivinar su clave secreta probando números.
    \item Copiar una firma vieja de esa persona.
    \item Pedirle que firme en secreto.
    \item Invertir el orden: con el desafío $c$ ya dado, escoger la respuesta $r$ y
      despejar $R$.
  \end{opciones}

\item Si simular no requiere la clave, ¿por qué no sirve para falsificar una firma de
  una sola persona?
  \begin{opciones}
    \item Porque el desafío debe salir del hash de $R$, y al simular se escoge $c$ antes
      de conocer $R$: casi nunca coinciden.
    \item Porque la simulación solo funciona con números pares.
    \item Porque el verificador pregunta al dueño de la clave.
    \item Sí sirve: cualquier firma se puede falsificar así.
  \end{opciones}

\item ¿En qué momento usa quien firma su clave secreta?
  \begin{opciones}
    \item Al principio, para escoger a los miembros del anillo.
    \item En cada uno de los eslabones.
    \item Nunca; por eso la firma es anónima.
    \item Al final, para cerrar el único hueco que queda: su propio eslabón.
  \end{opciones}

\item ¿Qué hace quien verifica para aceptar la firma?
  \begin{opciones}
    \item Recorre el anillo desde $c_1$ recalculando cada eslabón y comprueba que
      vuelve a obtener $c_1$.
    \item Pregunta a cada miembro si fue quien firmó.
    \item Comprueba que el firmante esté en una lista oficial.
    \item Usa la clave secreta del firmante para descifrar la firma.
  \end{opciones}

\item Alguien cambia una sola letra del mensaje después de que se firmó. ¿Qué pasa?
  \begin{opciones}
    \item Nada: la firma protege a los miembros, no al mensaje.
    \item La firma sigue siendo válida, pero cambia quién aparece como firmante.
    \item Los desafíos recalculados cambian, la vuelta ya no regresa a $c_1$ y la firma
      es inválida.
    \item Queda revelada la identidad de quien firmó.
  \end{opciones}

\item Una investigadora con una computadora infinitamente potente estudia una firma
  de anillo de 5 miembros. ¿Puede saber quién firmó?
  \begin{opciones}
    \item Sí, probando todas las claves hasta encontrar la correcta.
    \item Solo si el anillo tiene menos de 10 miembros.
    \item No: para cada miembro existe un camino válido que produce exactamente esa
      firma; lo mejor que puede hacer es adivinar (1 de 5).
    \item Sí, porque la clave del firmante siempre va primero.
  \end{opciones}

\item ¿Por qué Monero usa firmas de anillo enlazables?
  \begin{opciones}
    \item Para que la firma ocupe menos espacio.
    \item Para que un administrador pueda rastrear a los usuarios.
    \item Para detectar, con la imagen de clave, si la misma moneda se gasta dos veces,
      sin revelar quién la gasta.
    \item Para vincular cada transacción con el nombre real del usuario.
  \end{opciones}

\item En Monero cada anillo tiene 16 miembros. ¿Qué afirmación es correcta?
  \begin{opciones}
    \item La firma mide lo mismo sin importar el tamaño del anillo.
    \item La firma revela a la mitad de los miembros.
    \item Más miembros hacen que la firma sea más fácil de falsificar.
    \item Más miembros dan más anonimato (1 de 16 al adivinar), pero la firma crece con
      el anillo.
  \end{opciones}

\item ¿Por qué los ejemplos de este apunte NO son seguros?
  \begin{opciones}
    \item Porque usan números diminutos ($q = 11$, $23$ o $29$) y un hash de juguete, así
      que se pueden probar todas las opciones.
    \item Porque se calculan a mano o en un navegador.
    \item Porque las firmas de anillo están rotas en general.
    \item Porque el anillo tiene menos de 100 miembros.
  \end{opciones}
\end{enumerate}

\subsection*{Ejercicios para hacer a mano}
\addcontentsline{toc}{subsection}{Ejercicios para hacer a mano}

Usa el grupo $p = 23$, $q = 11$, $g = 4$ y la tabla de potencias de la
sección~\ref{sec:claves}.

\begin{enumerate}[leftmargin=2.4em,itemsep=0.5em]
  \item Beto tiene clave secreta $x = 5$. Calcula su clave pública y compárala con la
    del laboratorio 2.
  \item Con la clave de Ana ($x = 7$, $y = 8$) y $\mh = 0$, firma con $a = 9$: calcula
    $R$, $c$ y $r$, y comprueba la condición \ding{172}.
  \item En el laboratorio 2, repite la construcción suponiendo que firma
    \textbf{Diego} ($x = 9$) con $a = 5$ y que las respuestas simuladas son
    $r_1 = 2$, $r_2 = 10$, $r_3 = 5$, $r_5 = 8$. ¿Qué $r_4$ obtienes? Compara la firma
    con la de Carla y explica el resultado usando la sección~\ref{sec:quien}.
  \item Verifica la firma $\sigma = (1;\ 2, 10, 5, 5, 9)$ (solo cambió $r_5$) con el
    mensaje original. ¿Cierra el anillo? ¿En qué eslabón se nota el cambio?
\end{enumerate}

\clearpage
% ===========================================================================
\appendix
\section{Soluciones de la autoevaluación}\label{app:soluciones}
% ===========================================================================

\begin{center}\small
\captionsetup{type=table}
\caption{Respuestas correctas.}
\begin{tabular}{@{}*{12}{c}@{}}
\toprule
1 & 2 & 3 & 4 & 5 & 6 & 7 & 8 & 9 & 10 & 11 & 12 \\
\midrule
(b) & (c) & (a) & (d) & (a) & (d) & (a) & (c) & (c) & (c) & (d) & (a) \\
\bottomrule
\end{tabular}
\end{center}

\begin{enumerate}[leftmargin=2em,itemsep=0.6em,label=\textbf{\arabic*.}]
\item \textbf{(b).} Prueba que el autor está en el grupo y que el mensaje no se
  alteró. No dice quién fue y tampoco garantiza que el contenido sea cierto.
  \emph{Repasa la sección~\ref{sec:idea}.}
\item \textbf{(c).} Las claves públicas de los demás están a la vista de todos, así
  que no hace falta su permiso. Su propia clave secreta la usa una sola vez, para
  cerrar el anillo. \emph{Repasa la sección~\ref{sec:idea}.}
\item \textbf{(a).} La firma de grupo necesita un administrador que inscribe a los
  miembros y puede destapar al firmante. En la de anillo no hay administrador y el
  secreto es para siempre. \emph{Repasa la sección~\ref{sec:blockchain}.}
\item \textbf{(d).} Si escoges $r$ y despejas $R = g^{r} y^{c}$, la condición
  \ding{172} se cumple sin conocer la clave secreta. Es aritmética, no hay ningún
  secreto involucrado. \emph{Repasa la sección~\ref{sec:simular}.}
\item \textbf{(a).} La condición \ding{173} exige $c = H(m, R)$. Al simular eliges $c$
  primero y el hash de $R$ sale otro número. Solo acertarías por suerte, 1 vez de
  cada $q$, y en la práctica $q$ tiene 77 dígitos. \emph{Repasa la
  sección~\ref{sec:simular}.}
\item \textbf{(d).} Arranca en su eslabón con un número al azar, simula a todos los
  demás y, cuando la vuelta le regresa un desafío que no escogió, calcula
  $r = a - x c$. Solo ese paso exige la clave. \emph{Repasa la
  sección~\ref{sec:anillo}.}
\item \textbf{(a).} El verificador solo usa datos públicos: el mensaje, las claves
  públicas y $\sigma$. Si la vuelta regresa a $c_1$, la firma es válida. Además trata
  a todos los miembros igual. \emph{Repasa la sección~\ref{sec:verificar}.}
\item \textbf{(c).} La huella del mensaje entra en cada hash, así que cambian todos
  los desafíos y el anillo deja de cerrar. Con números de juguete puede coincidir por
  suerte; con números reales, no. \emph{Repasa la sección~\ref{sec:verificar}.}
\item \textbf{(c).} El anonimato es incondicional: la información sobre quién firmó no
  está en la firma. Ni conociendo todas las claves secretas se distingue al firmante.
  \emph{Repasa la sección~\ref{sec:quien}.}
\item \textbf{(c).} La imagen de clave es siempre la misma para la misma moneda, pero
  no indica cuál miembro del anillo es. Si aparece dos veces, la red rechaza el
  segundo gasto. \emph{Repasa la sección~\ref{sec:blockchain}.}
\item \textbf{(d).} En la versión de esta lección la firma tiene $n + 1$ números. Más
  miembros significa más señuelos entre los cuales esconderse, a cambio de una firma
  más grande y una verificación más lenta. \emph{Repasa la sección~\ref{sec:anillo}.}
\item \textbf{(a).} Con $q = 11$ basta probar 11 valores para encontrar cualquier
  clave secreta. Los sistemas reales usan números de unos 256 bits, normalmente sobre
  curvas elípticas, y un hash criptográfico. \emph{Repasa la sección~\ref{sec:claves}.}
\end{enumerate}

\begin{nota}{¿Cómo te fue?}
\textbf{12 aciertos:} entiendes qué promete una firma de anillo, cómo se construye y
por qué nadie puede saber quién firmó. \textbf{9 a 11:} muy bien; revisa la
explicación de las que fallaste. \textbf{6 a 8:} vas por buen camino; los
laboratorios 1 y 2 ayudan mucho. \textbf{Menos de 6:} conviene repasar desde la
sección~\ref{sec:simular}, la idea en la que se apoya todo lo demás.
\end{nota}

\section*{Referencias}
\addcontentsline{toc}{section}{Referencias}

\begin{itemize}[leftmargin=1.4em,itemsep=0.25em]
  \item R. Rivest, A. Shamir, Y. Tauman. \emph{How to Leak a Secret}. ASIACRYPT 2001.
  \item M. Abe, M. Ohkubo, K. Suzuki. \emph{1-out-of-n Signatures from a Variety of Keys}.
    ASIACRYPT 2002.
  \item C.-P. Schnorr. \emph{Efficient Identification and Signatures for Smart Cards}.
    CRYPTO 1989.
  \item J. K. Liu, V. K. Wei, D. S. Wong. \emph{Linkable Spontaneous Anonymous Group
    Signature for Ad Hoc Groups}. ACISP 2004.
  \item N. van Saberhagen. \emph{CryptoNote v2.0}. 2013.
  \item A. Kumar, C. Fischer, S. Tople, P. Saxena. \emph{A Traceability Analysis of
    Monero's Blockchain}. ESORICS 2017.
  \item M. Möser et al. \emph{An Empirical Analysis of Traceability in the Monero
    Blockchain}. PoPETs 2018.
  \item Fuente temática: \url{http://math.com.mx/anillo_de_firmas.html}.
\end{itemize}

\vfill
\begin{center}
\color{inkc}\small
\rule{0.35\textwidth}{0.4pt}\\[0.6em]
\textsc{ChainNotes} · Módulo 08 · Anterior: prueba de trabajo\\
{\footnotesize Versión LaTeX del apunte interactivo \texttt{temas/anillo-de-firmas.html}.
Todas las figuras se generan con TikZ.}
\end{center}

\end{document}
